\subsection{Matrix Groups}

Let $F$ be a field and let $n \in \zp$.

\begin{exercise}[1.4.1]
    Prove that $|\gl_2(\f_2)| = 6$.
\end{exercise}

\begin{sol}
    Direct calculation gives:
    \[\gl_2(\f_2) = 
    \set*{
    \begin{pmatrix}
        1 & 0 \\
        0 & 1
    \end{pmatrix}, 
    \begin{pmatrix}
        1 & 1 \\
        0 & 1
    \end{pmatrix},
    \begin{pmatrix}
        1 & 0 \\
        1 & 1
    \end{pmatrix},
    \begin{pmatrix}
        0 & 1 \\
        1 & 0
    \end{pmatrix},
    \begin{pmatrix}
        1 & 1 \\
        1 & 0
    \end{pmatrix},
    \begin{pmatrix}
        0 & 1 \\
        1 & 1
    \end{pmatrix}}. \qh
    \]
\end{sol}

\begin{exercise}[1.4.2]
    Write out all the elements of $\gl_2(\f_2)$ and compute the order of each element.
\end{exercise}

\begin{sol}
    By direct calculation, we obtain the following:
    \[
    \begin{array}{c|c|c|c|c|c|c}
        A & 
        \begin{pmatrix}
            1 & 0 \\
            0 & 1
        \end{pmatrix} & 
        \begin{pmatrix}
            1 & 1 \\
            0 & 1
        \end{pmatrix} &
        \begin{pmatrix}
            1 & 0 \\
            1 & 1
        \end{pmatrix} &
        \begin{pmatrix}
            0 & 1 \\
            1 & 0
        \end{pmatrix} &
        \begin{pmatrix}
            1 & 1 \\
            1 & 0
        \end{pmatrix} &
        \begin{pmatrix}
            0 & 1 \\
            1 & 1
        \end{pmatrix} \\
        \hline
        |A| & 1 & 2 & 2 & 2 & 3 & 3
    \end{array} \qh
    \]
\end{sol}

\begin{exercise} [1.4.3]
    Show that $\gl_2(\f_2)$ is non-Abelian.
\end{exercise}

\begin{sol}
    It is trivial to see that the following two matrices do not commute, hence $\gl_2(\f_2)$ is non-Abelian.
    \[
    \begin{pmatrix}
        1 & 1 \\
        0 & 1
    \end{pmatrix}, \quad
    \begin{pmatrix}
        1 & 1 \\
        1 & 0
    \end{pmatrix}. \qh
    \]
\end{sol}

\begin{exercise}[1.4.4]
    Show that if $n$ is not prime then $\intmod$ is not a field.
\end{exercise}

\begin{sol}
    If $n = 1$, then $\intmod[1] = \set{\bar 0}$ is not a field, since $\units[1] = \varnothing$, contradicting that $\units[1]$ should be a group. Otherwise, suppose $n \geq 2$ and is composite. Then there exist $a, b \in \zp$ such that $1 < a, b < n$ and $n = ab$. Moreover, $(a, n) \neq 1$ since $a \mid n$. By \exref{0.3.12}, there does not exist $c \in \z$ such that $ac \equiv 1 \bmod n$. Hence, $\intmod$ is not a field.
\end{sol}

\begin{exercise}[1.4.5]
    Show that $\gl_n(F)$ is a finite group if and only if $F$ has a finite number of elements.
\end{exercise}

\begin{solitem}
    \item [\rightimp] If $\gl_n(F)$ is finite, the subset $\set{\alpha I \mid \alpha\in F\unt}$ must also be finite, which implies that $F$ has a finite number of elements.
    \item [\leftimp] Suppose $|F| = q$. Then an $n \times n$ matrix over $F$ has $q^{n^2}$ possible entries. Since $\gl_n(F)$ is a subset of the set of all $n \times n$ matrices over $F$, it follows that $\gl_n(F)$ is finite.
\end{solitem}

\begin{exercise}[1.4.6]
    If $|F| = q$ is finite, prove that $|\gl_n(F)| < q^{n^2}$.
\end{exercise}

\begin{sol}
    See \leftimp in \exref{1.4.5}.
\end{sol}

\begin{exercise}[1.4.7]
    Let $p$ be a prime. Prove that the order of $\gl_2(\f_p)$ is $p^4 - p^3 - p^2 + p$.
    
    \note{Do not just quote the order formula in this section.}

    \hint{Subtract the number of $2 \times 2$ matrices which are \emph{not} invertible from the total number of $2 \times 2$ matrices over $\f_p$. You may use the fact that a $2 \times 2$ matrix is not invertible if and only if one row is a multiple of the other.}
\end{exercise}

\begin{sol}
    Observe that there are $p^4$ total $2 \times 2$ matrices over $\f_p$. The number of non-invertible matrices is as follows:
    \begin{itemize}
        \item If the first row is $(0, 0)$, then there are $p^2$ choices for the second row.
        \item If the first row is $(a, b) \neq (0, 0)$, then there are $p$ choices for the second row, since it must be a multiple of the first row. There are $p^2 - 1$ choices for the first row.
    \end{itemize}
    Then the number of non-invertible matrices is $p^2 + p(p^2 - 1) = p^3 + p^2 - p$. Subtracting from the total number of matrices, we have the desired number.
\end{sol}

\begin{exercise}[1.4.8]
    Show that $\gl_n(F)$ is non-Abelian for any $n \geq 2$ and any $F$.
\end{exercise}

\begin{sol}
    Consider the subset $S = \set{0, 1}$ of $F$. Then $\gl_n(S)$ is a subset of $\gl_n(F)$. We proceed by induction on this subset. For $n = 2$, we have $\gl_2(S)$ is non-Abelian by \exref{1.4.3}. Suppose that $\gl_n(S)$ is non-Abelian for some $n \geq 2$. Then there exist $A, B \in \gl_n(S)$ such that $AB \neq BA$. Consider the matrices
    \[A_0 = 
    \begin{pmatrix}
        A & 0 \\
        0 & 1
    \end{pmatrix},
    B_0 =
    \begin{pmatrix}
        B & 0 \\
        0 & 1
    \end{pmatrix} \in \gl_{n + 1}(S).\]
    Observe that $A_0B_0$ contains $AB$ as a block matrix, while $B_0A_0$ contains $BA$ as a block matrix. Since $AB \neq BA$, then $A_0B_0 \neq B_0A_0$. Hence, $\gl_{n + 1}(S)$ is non-Abelian. By induction, $\gl_n(S)$ is non-Abelian for all $n \geq 2$. Since $\gl_n(S) \subseteq \gl_n(F)$, then $\gl_n(F)$ is non-Abelian for all $n \geq 2$.
\end{sol}

\begin{exercise}[1.4.9]
    Prove that the binary operation of matrix multiplication of $2 \times 2$ matrices with real number entries is associative.
\end{exercise}

\begin{sol}
    The direct proof is laborious. Instead, associate with any $2 \times 2$ matrix some linear map. Since composition of linear maps is associative, the result follows.
\end{sol}

\begin{exercise}[1.4.10]
    Let
    \[G = \left\{
    \begin{pmatrix}
        a & b \\
        0 & c
    \end{pmatrix}~\middle|~a, b, c \in \r, a \neq 0, c \neq 0 \right\}\]
    \begin{subexercises}
        \item Compute the product of
        \[
        \begin{pmatrix}
            a_1 & b_1 \\
            0 & c_1
        \end{pmatrix} \text{ and }
        \begin{pmatrix}
            a_2 & b_2 \\
            0 & c_2
        \end{pmatrix}
        \]
        to show that $G$ is closed under matrix multiplication.
        \item Find the matrix inverse of
        \[
        \begin{pmatrix}
            a & b \\
            0 & c
        \end{pmatrix}\]
        and deduce that $G$ is closed under inverses.
        \item Deduce that $G$ is a \defref{subgroup} of $\gl_2(\r)$.
        \item Prove that the set of elements of $G$ whose two diagonal entries are equal (i.e., $a = c$) is also a \defref{subgroup} of $\gl_2(\r)$.
    \end{subexercises}
\end{exercise}

\begin{solenum}
    \item 
    \[\begin{pmatrix}
        a_1 & b_1 \\
        0 & c_1
    \end{pmatrix}
    \begin{pmatrix}
        a_2 & b_2 \\
        0 & c_2
    \end{pmatrix} = 
    \begin{pmatrix}
        a_1a_2 & a_1b_2 + b_1c_2 \\
        0 & c_1c_2
    \end{pmatrix}\]
    Since $a_i \neq 0$ and $c_i \neq 0$, then $a_1a_2 \neq 0$ and $c_1c_2 \neq 0$ so that $G$ is closed under matrix multiplication.
    \item The determinant is $ac$. We then have the inverse
    \[
    \begin{pmatrix}
        a & b \\
        0 & c
    \end{pmatrix}\inv = 
    \frac{1}{ac}
    \begin{pmatrix}
        c & -b \\
        0 & a
    \end{pmatrix} = 
    \begin{pmatrix}
        1/a & -b/ac \\
        0 & 1/c
    \end{pmatrix}\]
    \item Since $G$ is closed under inverses and matrix multiplication, the result follows.
    \item Let $H = \set{A \in G \mid a = c}$. Then
    \[
    \begin{pmatrix}
        a_1 & b_1 \\
        0 & a_1
    \end{pmatrix}
    \begin{pmatrix}
        a_2 & b_2 \\
        0 & a_2
    \end{pmatrix} = 
    \begin{pmatrix}
        a_1a_2 & a_1b_2 + b_1a_2 \\
        0 & a_1a_2
    \end{pmatrix} \in H
    \]
    and
    \[
    \begin{pmatrix}
        a & b \\
        0 & a
    \end{pmatrix}\inv = 
    \frac{1}{a^2}
    \begin{pmatrix}
        a & -b \\
        0 & a 
    \end{pmatrix} = 
    \begin{pmatrix}
        1/a & -b/a^2 \\
        0 & 1/a
    \end{pmatrix} \in H
    \]
    Then $H$ is closed under inverses and matrix multiplication, hence it is a subgroup of $\gl_2(\r)$.
\end{solenum}

\begin{exercise} [1.4.11]
    Let
    \[H(F) = \left\{\left.
    \begin{pmatrix}
        1 & a & b \\
        0 & 1 & c \\
        0 & 0 & 1
    \end{pmatrix}\,\right|\, a, b, c \in F\right\},\]
    called the \textbf{Heisenberg group} over $F$. Let
    \[X = 
    \begin{pmatrix}
        1 & a & b \\
        0 & 1 & c \\
        0 & 0 & 1
    \end{pmatrix} \text{ and }Y = 
    \begin{pmatrix}
        1 & d & e \\
        0 & 1 & f \\
        0 & 0 & 1
    \end{pmatrix}\]
    be elements of $H(F)$.
    \begin{subexercises}
        \item Compute the matrix product $XY$ and deduce that $H(F)$ is closed under matrix multiplication. Exhibit explicit matrices such that $XY \neq YX$ (so that $H(F)$ is always non-Abelian).
        \item Find an explicit formula for the matrix inverse $X\inv$ and deduce that $H(F)$ is closed under inverses.
        \item Prove the associative law for $H(F)$ and deduce that $H(F)$ is a group of order $|F|^3$. (Do not assume that matrix multiplication is associative.)
        \item Find the order of each element of the finite group $H(\intmod[2])$.
        \item Prove that every non-identity of the group $H(\r)$ has infinite order.
    \end{subexercises}
    \mnote{The choice of rewriting the matrix as a tuple is simply a time saver/ease of notation for the sake of computation. However, it does allow one to encode $H(F)$ in the multiplication law of the field.}
\end{exercise}

\begin{solenum}
    \item Calculating the product, we have
    \[XY = 
    \begin{pmatrix}
        1 & d + a & e + af + b \\
        0 & 1 & f + c \\
        0 & 0 & 1
    \end{pmatrix} \in H(F)\]
    Consider the matrices
    \[X =
    \begin{pmatrix}
        1 & 1 & 1 \\
        0 & 1 & 0 \\
        0 & 0 & 1
    \end{pmatrix}, \quad
    Y =
    \begin{pmatrix}
        1 & 0 & 0 \\
        0 & 1 & 1 \\
        0 & 0 & 1
    \end{pmatrix} \in H(F).\]
    Then
    \[XY =
    \begin{pmatrix}
        1 & 1 & 1 \\
        0 & 1 & 1 \\
        0 & 0 & 1
    \end{pmatrix} \neq
    \begin{pmatrix}
        1 & 1 & 0 \\
        0 & 1 & 1 \\
        0 & 0 & 1
    \end{pmatrix} = YX,\]
    so that $H(F)$ is non-Abelian.
    \item From part (a), we know $G$ is closed under matrix multiplication. For ease of notation, represent a matrix as a tuple: $X = (a, b, c)$ and $Y = (d, e, f)$. Then we may represent matrix multiplication in $H(F)$ as
    \[(a, b, c)(d, e, f) = (a + d, af + b + e, c + f)\]
    Now we want to find $(x, y, z)$ such that $(a, b, c)(x, y, z) = (0, 0, 0)$. Then $(a + x, az + b + y, c + z) = (0, 0, 0)$ implies $x = -a$, $y = -b + ac$, and $z = -c$. It follows that
    \[X\inv = 
    \begin{pmatrix}
        1 & -a & -b + ac \\
        0 & 1 & -c \\
        0 & 0 & 1
    \end{pmatrix} = (-a, -b + ac, -c)\]
    so that $H(F)$ is closed under inverses.
    \item For some $(a, b, c) \in H(F)$, we know that there are $|F|$ choices for each of $a, b, c$. Then $|H(F)| = |F|^3$. Now let $(a, b, c), (d, e, f), (g, h, i) \in H(F)$. Then we have
    \begin{align*}
        ((a, b, c)(d, e, f))(g, h, i) & = (a + d, af + b + e, c + f)(g, h, i) \\
        & = (a + d + g, (a + d)i + af + b + e + h, c + f + i) \\
        & = (a + d + g, ai + di + af + b + e + h, c + f + i) \\
        & = (a, b, c)(d + g, e + h, f + i) \\
        & = (a, b, c)((d, e, f)(g, h, i))
    \end{align*}
    \item The elements of $H(\intmod[2])$ are
    \[ (0, 0, 0), (1, 0, 0), (0, 1, 0), (0, 0, 1), (1, 1, 0), (1, 0, 1), (0, 1, 1), (1, 1, 1).\]
    Since $(a, b, c)^2 = (2a, 2b + ac, 2c) = (0, ac, 0)$ in $\intmod[2]$, then we have two cases:
    \begin{itemize}
        \item If $ac = 0$, then at least one of $a$ or $c$ is 0, while $b$ varies. These are the matrices that have order 2:
        \[(1, 0, 0), (0, 1, 0), (0, 0, 1), (1, 1, 0), (0, 1, 1).\]
        \item If $ac = 1$, then both $a$ and $c$ are 1, while $b$ varies. These are the matrices that have order 4:
        \[(1, 0, 1), (1, 1, 1).\]
    \end{itemize}
    \item Let $X \in H(\r)$ be non-identity. It is easy to see by induction that
    \[X^n = (na, n(n - 1)ac/2 + nb, nc).\]
    Since $X \neq I_3$, then at least one of $a, b, c$ is nonzero. It is clear that $a \neq 0$ or $c \neq 0$ implies $X^n \neq I_3$. Now suppose $a = c = 0$ and $b \neq 0$. Then $X^n = (0, nb, 0)$. Since $b \neq 0$, then $nb \neq 0$ for all $n \in \zp$. Hence, $X^n \neq I_3$ for all $n \in \zp$. Therefore, every non-identity element of $H(\r)$ has infinite order.
\end{solenum}